\documentclass[11pt,a4paper]{article}
\usepackage{fontspec}
\usepackage[french]{babel}
\usepackage{amsmath,amssymb}
\usepackage[margin=2.3cm]{geometry}
\usepackage{booktabs,longtable,tabularx,array}
\usepackage{enumitem}
\usepackage{xcolor}
\usepackage{microtype}
\usepackage{url}
\usepackage{hyperref}

\hypersetup{colorlinks=true,linkcolor=blue!45!black,urlcolor=blue!45!black}
\setlist{nosep}
\newcommand{\code}[1]{\texttt{#1}}
\newcommand{\NW}{\mathrm{NW}}
\newcommand{\E}{\mathbb{E}}

\title{Protocole de l'étude de sensibilité du modèle M4B\\
\large Cartographie des paramètres, plan d'expérience et règles de confirmation}
\author{Campagne \code{recherche/sensibilite\_m4b}}
\date{17 juillet 2026}

\begin{document}
\maketitle

\begin{abstract}
Ce protocole fige, avant les simulations de production, la conduite d'une
étude de sensibilité du modèle de société de crédit M4B. Il définit l'objet
étudié, le statut expérimental de chaque paramètre, le dictionnaire des
grandeurs de réponse, le plan d'expérience en six volets (contrôle de
taille, contrôle temporel, courbes un-facteur-à-la-fois, plan global
espace-remplissant, coupes bidimensionnelles, phase confirmatoire), le
diagnostic obligatoire du garde-fou de population et les règles de passage
de l'exploration à la confirmation. Les plages initiales sont des points de
départ soumis à un audit par cellules pilotes ; leur révision éventuelle
sera consignée dans le rapport préliminaire. L'audit du moteur (tests
unitaires et comptables, parité avec le modèle source M4, neutralité des
options de mesure) et le chiffrage du coût de calcul sont rapportés ici car
ils conditionnent le budget du plan.
\end{abstract}

\tableofcontents

\section{Objet, périmètre et principes}

\subsection{Objet}

Le modèle étudié est M4B, moteur minimal de société de crédit installé
dans le dossier \path{m4b_credit_soc_mini} du dépôt. M4B reproduit la
mécanique validée du modèle M4 (dossier \path{m4_credit_soc_fable}, traité
en lecture seule comme référence historique) en supprimant toutes
les variantes d'ablation. L'étude doit cartographier honnêtement :

\begin{enumerate}
\item quels paramètres contrôlent le régime démographique, le réseau de
  crédit, les faillites en cascade et les distributions individuelles ;
\item où se trouvent les non-linéarités, seuils, interactions et
  changements de régime ;
\item quelles signatures sont robustes aux paramètres, aux graines, à la
  durée et aux choix de fenêtre statistique ;
\item quels paramètres semblent actifs, redondants ou assimilables à des
  conventions d'échelle, afin d'éclairer une éventuelle réduction du
  modèle avant son successeur M5.
\end{enumerate}

L'objectif n'est ni de chercher la configuration qui « donne le plus de
criticité auto-organisée », ni de ne retenir que les résultats compatibles
avec les rapports antérieurs. Les extinctions, explosions, absences d'effet
et contradictions sont des résultats à conserver et à expliquer.

\subsection{Recours exceptionnel à M4}

Aucune campagne M4 n'est menée. Le seul recours exécutable à M4 admis par
ce protocole est le test de parité existant
(\code{tests/test\_mini.py}), exécuté une fois comme contrôle de provenance
de la réduction (section~\ref{sec:audit}). Toute affirmation de la présente
étude est vérifiée sur M4B.

\section{Le modèle M4B}
\label{sec:modele}

Cette section résume la mécanique exécutée, d'après le code du moteur
(\code{m4b/model.py}) et son rapport de mécanique
(\code{report/mecanique\_m4b.pdf}). Elle rend le protocole autonome ; le
rapport de mécanique reste la référence détaillée.

\subsection{État}

Chaque entité $i$ possède un unique stock réel, son capital productif
$K_i(t)\ge 0$, exprimé en joules (J), l'unité de ressource du projet. Les
créances et dettes proviennent exclusivement de contrats de prêt
$(\ell,b,q,r)$ : prêteuse $\ell$, emprunteuse $b$, principal nominal $q>0$
(J), taux $r>0$ par pas. Les contrats sont perpétuels : seuls les intérêts
$rq$ sont servis, le principal n'est jamais amorti et ne se déprécie pas.
Deux prêts successifs entre la même paire orientée fusionnent au taux moyen
pondéré. Pour chaque entité, $C_i$ est le total de ses créances, $D_i$ sa
dette nominale, et sa valeur nette vaut
\[
\NW_i = K_i + C_i - D_i .
\]

\subsection{Chronologie d'un pas}

L'ordre des huit phases fait partie du modèle :

\begin{enumerate}
\item \textbf{Naissances} : $B_t\sim\mathrm{Poisson}(\lambda)$ entités
  naissent avec $K_0$ et aucun contrat.
\item \textbf{Choc} : $K_i \leftarrow K_i e^{\eta_i}$,
  $\eta_i \overset{\text{iid}}{\sim}
  \mathcal N(-\sigma^2/2,\,\sigma^2)$, donc $\E[e^{\eta_i}]=1$.
\item \textbf{Production} : $K_i \leftarrow K_i + \alpha\sqrt{K_i}$, avec
  $\alpha=1$ (convention d'échelle, hors campagne).
\item \textbf{Intérêts} : chaque emprunteuse paie son service
  $S_b=\sum r q$ si $K_b\ge S_b$ ; sinon elle verse tout son capital au
  prorata et fait \emph{défaut de service}.
\item \textbf{Dépréciation} : $K_i \leftarrow (1-\delta)K_i$ ; les
  principaux nominaux sont inchangés.
\item \textbf{Marché} : un round par entité vivante non défaillante ; on
  tire $k$ entités distinctes, la plus riche prête à la plus pauvre au taux
  $r=\sqrt{r_\ell r_b}$ avec $r_x = \alpha/(2\sqrt{K_x})$, vers la cible
  commune $K^\ast=\sqrt{K_\ell K_b}$ ; le transfert
  $q=\min(K_\ell-K^\ast,\,\max(0,K^\ast-K_b))$ conserve le capital total et
  les valeurs nettes des deux parties.
\item \textbf{Faillites} : toute vivante en défaut de service ou de valeur
  nette négative meurt ; ses prêteuses perdent l'intégralité de leur
  créance, ses propres créances sont annulées et son capital résiduel est
  détruit ; les survivantes devenues insolvables forment la génération
  suivante, jusqu'au point fixe (cascade).
\item \textbf{Mesure} : séries, avalanches, instantanés ; aucune phase de
  mesure ne consomme de tirage aléatoire.
\end{enumerate}

Une avalanche est une composante faiblement connexe du graphe des arêtes de
perte restreint aux mortes du pas ; sa taille $s$ est son nombre de mortes,
sa profondeur la génération maximale, ses racines les membres de la
première file. La mesure est strictement intra-pas.

\subsection{Statuts d'arrêt}

Une exécution s'arrête à l'horizon $T$, ou avant si la population atteint
zéro (statut \code{extinction}, absorbant) ou dépasse \code{pop\_max}
(statut \code{explosion} dans le vocabulaire du moteur). Ce protocole ne
reprend pas ce vocabulaire à son compte : atteindre \code{pop\_max} est une
\emph{censure numérique} qui déclenche le diagnostic de la
section~\ref{sec:popmax}, jamais une conclusion.

\section{Statut expérimental des paramètres}
\label{sec:parametres}

\begin{center}
\begin{tabularx}{\textwidth}{@{}llX@{}}
\toprule
Paramètre & Défaut & Statut dans l'étude \\
\midrule
$\lambda$ (\code{lam}) & 10 & scientifique ; traité d'abord comme axe de
taille finie et démographique (contrôle de scaling), ensuite seulement
comme levier économique \\
$\delta$ (\code{delta}) & 0{,}05 & scientifique : dépréciation du capital \\
$\sigma$ (\code{sigma}) & 0{,}25 & scientifique : volatilité du choc \\
$K_0$ (\code{K0}) & 25 & scientifique : capital de naissance \\
$k$ (\code{k}) & 3 & scientifique, discret : taille de l'échantillon du
marché local \\
\code{seed} & --- & réplication stochastique, jamais un mécanisme \\
$T$ (\code{T}) & 2000 & durée d'observation : contrôle de probité
temporelle uniquement (section~\ref{sec:temps}) \\
\code{pop\_max} & 30\,000 & garde-fou d'arrêt : censure, jamais une borne
dynamique \\
$\alpha=1$, tolérances & --- & conventions et constantes du moteur, non
promues en paramètres \\
\bottomrule
\end{tabularx}
\end{center}

Le classement d'importance final portera sur $\delta,\sigma,K_0,k$ (leviers
microscopiques) d'une part, et sur $\lambda$ (taille) d'autre part, sans
les fusionner dans un classement unique non qualifié. $T$, le burn-in et
les fenêtres n'entrent dans aucun classement.

\subsection{Rôle strict de l'étude temporelle}
\label{sec:temps}

La variation de $T$, du burn-in et des fenêtres sert exclusivement à
vérifier la probité des résultats. Une conclusion n'est robuste dans le
temps que si : elle persiste lorsque l'horizon est prolongé ; elle ne
dépend pas d'une seule définition du burn-in ; les métriques cumulatives et
par fenêtres donnent une lecture compatible ; les dérives résiduelles sont
quantifiées et visibles. Le burn-in par défaut est $T/4$ ; les burn-ins
$T/8$ et $T/2$ ainsi que la coupure de la fenêtre par défaut en deux
demi-fenêtres sont systématiquement calculés. Pour les régimes lents
(petits $\sigma$ et $\delta$), $T$ pourra être poussé au-delà de 10\,000,
avec un suivi explicite du coût de calcul et de stockage.

\section{Questions scientifiques}

\begin{enumerate}
\item \textbf{Régime} : dans quelles régions de l'espace
  $(\delta,\sigma,K_0,k)$ la population est-elle stationnaire, éteinte, ou
  en croissance soutenue ? Où sont les frontières et de quelle nature
  sont-elles (brutales ou graduelles) ?
\item \textbf{Taille} : quelles grandeurs sont intensives (indépendantes de
  $\lambda$ après normalisation) et lesquelles portent un effet de taille
  finie (coupure des avalanches, maximum absolu) ?
\item \textbf{Crédit et cascades} : comment l'exposition (dette/capital,
  degrés, concentration) et la propagation (rapport de branchement,
  fraction induite, profondeur) répondent-elles aux paramètres ?
\item \textbf{Distributions} : comment les distributions individuelles
  ($K$, valeur nette, revenu, âge) et les inégalités (Gini, parts hautes)
  se déforment-elles ?
\item \textbf{Robustesse} : quelles signatures (exposant de taille
  d'avalanche, rapport de branchement, $N/\lambda$) sont stables sous
  graines, durées et fenêtres ?
\item \textbf{Réduction} : quels paramètres sont actifs, faiblement
  identifiables, redondants ou assimilables à des conventions d'échelle ?
\end{enumerate}

\section{Dictionnaire des métriques}
\label{sec:metriques}

Toutes les métriques fenêtrées sont calculées sur la fenêtre par défaut
$[T/4,\,t_{\mathrm{fin}}]$ (et sur les fenêtres de contrôle de la
section~\ref{sec:temps}), où $t_{\mathrm{fin}}\le T$ est le dernier pas
exécuté. Une fenêtre de moins de 30 pas est déclarée invalide et la cellule
est traitée comme non mesurable (mais conservée et comptée). L'extraction
est implémentée dans \code{scripts/lib\_metrics.py} ; chaque run produit un
JSON de métriques réutilisable.

\subsection{Régime et démographie}

\begin{longtable}{p{0.27\textwidth}p{0.40\textwidth}p{0.08\textwidth}p{0.15\textwidth}}
\toprule
Métrique & Définition & Unité & Statut \\
\midrule
\endhead
statut & \code{ok} / \code{extinction} / censure par \code{pop\_max} &
--- & primaire \\
$\bar N$ & moyenne de la population sur la fenêtre & entités & primaire \\
$\bar N/\lambda$ & taille normalisée par l'intensité des naissances
(loi de Little : $\bar N/\lambda$ estime l'espérance de durée de vie) &
pas & primaire \\
pente relative & pente OLS de $N(t)$ multipliée par la longueur de fenêtre,
rapportée à $\bar N$ ; erreur type corrigée par le temps d'autocorrélation
intégré $\tau_{\mathrm{int}}$ & --- & primaire \\
$\tau_{\mathrm{int}}$, ESS & temps d'autocorrélation intégré de $N(t)$ ;
taille d'échantillon effective $n/\tau_{\mathrm{int}}$ & pas & secondaire \\
CV de $N$ & écart-type/moyenne de $N(t)$ & --- & secondaire \\
naissances, morts & moyennes par pas ; taux de mortalité
$\overline{\mathrm{morts}}/\bar N$ & pas$^{-1}$ & secondaire \\
âge au décès & moyenne, médiane, q90 des âges dans \code{deaths.csv} &
pas & secondaire \\
stationnarité & écart relatif de $\bar N$ entre les deux demi-fenêtres, et
en unités d'écart-type regroupé & --- & primaire (validité) \\
\bottomrule
\end{longtable}

\subsection{Crédit, exposition et réseau}

\begin{longtable}{p{0.27\textwidth}p{0.40\textwidth}p{0.08\textwidth}p{0.15\textwidth}}
\toprule
Métrique & Définition & Unité & Statut \\
\midrule
\endhead
prêts par tête & $\overline{n_{\mathrm{prêts}}}/\bar N$ & --- & primaire \\
dette/capital & $\sum D_i / \sum K_i$ sur l'instantané final & --- &
primaire \\
crédit nouveau / capital & volume de prêts nouveaux par pas rapporté à
$\overline{K_{\mathrm{tot}}}$ & pas$^{-1}$ & secondaire \\
intérêts / capital & intérêts payés par pas rapportés à
$\overline{K_{\mathrm{tot}}}$ & pas$^{-1}$ & secondaire \\
pertes / capital & créances détruites par pas rapportées à
$\overline{K_{\mathrm{tot}}}$ & pas$^{-1}$ & secondaire \\
degrés & moyennes et maxima des degrés prêteur et emprunteur (instantané
final) & --- & secondaire \\
concentration & Gini des principaux ; part du 1\,\% des plus gros contrats
& --- & secondaire \\
fractions exposées & parts d'entités emprunteuses et prêteuses & --- &
secondaire \\
\bottomrule
\end{longtable}

\subsection{Avalanches et propagation causale}

Événements de \code{avalanches.csv} avec $t$ dans la fenêtre. Le seuil bas
des ajustements est $s_{\min}=2$ (les tailles 1 sont des morts isolées).

\begin{longtable}{p{0.27\textwidth}p{0.40\textwidth}p{0.08\textwidth}p{0.15\textwidth}}
\toprule
Métrique & Définition & Unité & Statut \\
\midrule
\endhead
fréquence & événements par pas & pas$^{-1}$ & secondaire \\
quantiles de taille & médiane, q90, q99, maximum ; maximum rapporté à
$\bar N$ & --- & primaire \\
susceptibilité & $\langle s^2\rangle/\langle s\rangle$ & --- & primaire \\
rapport de branchement & $b = 1 - \sum_a R_a / \sum_a S_a$ ($R_a$ racines,
$S_a$ taille) ; fraction induite $1-b$ complémentaire & --- & primaire \\
profondeur & moyenne et maximum des générations & --- & secondaire \\
exposant $\hat\alpha$ & MLE discrète de $p(s)=s^{-\alpha}/\zeta(\alpha,2)$,
$s\ge 2$ ; erreur type asymptotique & --- & primaire \\
loi tronquée & MLE discrète de $p(s)\propto s^{-\alpha}e^{-s/s_c}$ ;
coupure $s_c$ & --- & primaire \\
LRT coupure & $2(\ell_{\mathrm{tronquée}}-\ell_{\mathrm{pure}})$, p-valeur
conservatrice $\tfrac12\chi^2_1$ (paramètre au bord) & --- & secondaire \\
Vuong PL/LN & test de Vuong loi de puissance contre log-normale discrète
renormalisée ($z>0$ : puissance favorisée) & --- & primaire \\
volume causal & moyenne et maximum de \code{volume\_j} & J & secondaire \\
\bottomrule
\end{longtable}

\textbf{Règle d'identifiabilité.} Un ajustement de loi n'est tenté que si
la fenêtre contient au moins 100 avalanches de taille $\ge 2$ ; sinon la
métrique est marquée \emph{non identifiable} et aucun nombre n'est produit.
Jamais de conclusion de loi de puissance sur le seul $r^2$ d'une régression
log-log. Les événements de graines différentes ne sont \emph{jamais}
fusionnés avant un ajustement : chaque graine est ajustée séparément, puis
les estimateurs sont agrégés (moyenne, écart-type inter-graines).

\subsection{Distributions individuelles et inégalités}

Sur l'instantané final (et, pour les runs confirmatoires, sur les
instantanés périodiques après burn-in) : moyenne, médiane, q90, q99,
maximum, Gini (partie positive), part du décile supérieur et écart-type
logarithmique de $K$, $\NW$, revenu brut ($P_i + $ intérêts reçus) et
revenu net. Contrôles : corrélation de Spearman âge--$K$ (effet de
cohorte), part d'entités à valeur nette négative, renouvellement du haut de
la distribution (pour les runs à instantanés périodiques : recouvrement du
top 1\,\% entre instantanés successifs). Statut : Gini de $K$ et part du
décile sont primaires, le reste secondaire. Les familles sont comparées sur
les mêmes supports et avec les mêmes règles de troncature.

\subsection{Intégrité numérique}

Par run : identité comptable
$\Delta K_{\mathrm{tot}} = I_t + G_t + P_t - \Delta_t - X_t$ (résidu
relatif maximal), contrôle du carnet (\code{book\_errors} de
\code{summary.json} : contrats orphelins, agrégats, créances $=$ dettes),
$K\ge0$, durée d'exécution, taille disque, statut d'arrêt. Un run avec
résidu relatif $>10^{-9}$ ou erreurs de carnet est exclu des estimateurs et
compté comme anomalie dans le journal.

\section{Plan d'expérience}
\label{sec:plan}

\subsection{Référence et centre de campagne}

La référence du modèle est
\[
\lambda=10,\quad \delta=0{,}05,\quad \sigma=0{,}25,\quad K_0=25,\quad
k=3,\quad T=2000 .
\]
Le centre de la campagne est le même point à $\lambda=30$. Justification :
au benchmark (section~\ref{sec:audit}), la fenêtre par défaut d'un run de
référence contient $\sim$2\,100 avalanches de taille $\ge2$ contre
$\sim$5\,600 à $\lambda=30$, pour un coût de 80\,s par run ; les queues y
sont mieux identifiées sans changer les grandeurs intensives
($\bar N/\lambda = 14{,}8$ aux deux tailles). Le lien avec la référence est
maintenu par l'axe de taille $\lambda\in\{10,30,100\}$, mesuré avec le
même panel de graines.

\subsection{Panels de graines}

\begin{itemize}
\item \textbf{Pilotes} : graines $\{1,2\}$.
\item \textbf{Exploration} (OAT, plan global, coupes 2D) : graines
  $\{1,2,3\}$, identiques pour toutes les cellules d'une même comparaison
  (contrastes appariés).
\item \textbf{Confirmation} : graines $\{11,\dots,15\}$, disjointes de
  l'exploration ; augmentées au-delà de cinq près d'une frontière de
  régime ou lorsque la variance inter-graines domine l'effet.
\end{itemize}

Le partage d'un panel permet des contrastes appariés entre cellules ; on ne
parlera pas de « common random numbers » au sens strict, car deux
configurations consomment ensuite des tirages en nombre différent.

\subsection{Volets successifs}

\begin{enumerate}
\item \textbf{Pilotes d'audit des plages} (graines $\{1,2\}$, $T=2000$) :
  les coins de l'espace $(\delta,\sigma)\in\{0{,}02;0{,}10\}\times
  \{0{,}10;0{,}50\}$, les extrêmes $K_0\in\{5,100\}$ et $k\in\{2,10\}$, les
  contrôles $\sigma=0$ (avec $k=3$, $k=2$ et $\delta=0{,}02$), le coin de
  mortalité $(K_0{=}5,\sigma{=}0{,}50,\delta{=}0{,}10)$, le centre, un
  contrôle temporel $T=4000$ au centre, une sonde lente
  $(\delta{=}0{,}02,\sigma{=}0{,}10)$ à $T=8000$ et une sonde de coût
  $\lambda=100$ lente. Critères de révision des bornes : une plage n'est
  retenue que si elle produit au moins un régime non trivial ; un coin qui
  ne produit que des extinctions immédiates ou des censures est resserré ou
  documenté comme régime propre.
\item \textbf{Contrôle de taille} : $\lambda\in\{10,30,100\}$ au centre
  microscopique, graines d'exploration puis de confirmation ; lecture des
  grandeurs intensives et du scaling de la coupure $s_c$ et du maximum.
\item \textbf{Contrôle temporel} : $T\in\{2000,4000\}$ au centre et sur
  les cellules confirmatoires ; burn-ins $T/8$, $T/4$, $T/2$ ;
  demi-fenêtres. Prolongation au-delà de $T=8000$ uniquement pour les
  régimes lents avérés.
\item \textbf{Courbes OAT} autour du centre ($T=2000$, graines
  $\{1,2,3\}$) :
  $\delta\in\{0{,}02;0{,}03;0{,}05;0{,}07;0{,}10\}$,
  $\sigma\in\{0;0{,}10;0{,}15;0{,}20;0{,}25;0{,}325;0{,}40;0{,}50\}$,
  $K_0\in\{5;10;25;50;100\}$ (échelle logarithmique),
  $k\in\{2;3;4;6;10\}$, $\lambda\in\{10;30;100\}$. Objectif : monotonies,
  seuils, formes non linéaires, avec incertitude inter-graines et
  contrastes appariés.
\item \textbf{Plan global espace-remplissant} : hypercube latin stratifié
  de 64 points sur $(\delta,\sigma,K_0)$ continus
  ($\delta\in[0{,}02;0{,}10]$ et $\sigma\in[0{,}05;0{,}50]$ linéaires,
  $K_0\in[5;100]$ log-uniforme) avec équilibrage du facteur discret
  $k\in\{2,3,4,6,10\}$ (12 ou 13 points par niveau), $\lambda=30$,
  $T=2000$, graines $\{1,2,3\}$, soit 192 runs. Le plan est engendré par un
  générateur pseudo-aléatoire dédié de graine 20260717 et sauvegardé dans
  \code{manifests/}. Aucun indice de Sobol ne sera revendiqué : le plan
  n'est pas un plan de Saltelli.
\item \textbf{Coupes 2D ciblées} : deux à trois cartes $7\times7$ (ou
  résolution adaptée) choisies \emph{après} le screening sur les
  interactions et frontières candidates, graines $\{1,2,3\}$.
\item \textbf{Phase confirmatoire} : cellules retenues rejouées à
  $T=4000$, graines $\{11,\dots,15\}$, écriture enrichie (instantanés
  périodiques) ; robustesse aux fenêtres exigée. Les runs confirmatoires
  sont versés dans Simulation Lab avec figures et métadonnées, marqués à
  conserver.
\end{enumerate}

\subsection{Analyses de sensibilité}

Lectures complémentaires, jamais un classement unique sans métrique
précisée :

\begin{itemize}
\item courbes OAT avec incertitude inter-graines et contrastes appariés ;
\item corrélations de rang partielles (PRCC) et régressions linéaires
  standardisées sur le plan global, avec intervalles bootstrap (rééchantillonnage
  des points du plan) ;
\item surface de réponse non linéaire (polynôme de degré 2 avec
  interactions, moindres carrés) validée hors échantillon par validation
  croisée ; importance par permutation si le pouvoir prédictif hors
  échantillon le justifie ($R^2_{\mathrm{CV}}>0{,}5$ typiquement) ;
\item cartes de régime extinction / stationnaire / croissance-censure et
  diagrammes 2D des interactions confirmées ;
\item décomposition de la variabilité : pour chaque métrique, part de
  variance entre cellules (paramètres) et intra-cellule (graines), estimée
  sur le plan global.
\end{itemize}

Toute analyse exploratoire ayant servi à choisir des cellules est
étiquetée comme telle et séparée de la confirmation.

\section{Diagnostic obligatoire du garde-fou \code{pop\_max}}
\label{sec:popmax}

\code{pop\_max} censure une trajectoire par sécurité ; l'atteindre ne
démontre jamais une explosion. Pour toute cellule censurée :

\begin{enumerate}
\item répétition avec plafonds élevés (\code{pop\_max} $\in$
  \{30\,000 ; 60\,000 ; 120\,000\}) et le même panel de graines ;
\item examen des pentes, accélérations et temps de franchissement des
  plafonds successifs, plutôt que de la seule population finale ;
\item encadrement adaptatif de la frontière par des cellules voisines des
  deux côtés dans l'espace des paramètres ;
\item vérification d'un changement qualitatif entre un régime voisin
  stationnaire (qui peut frôler un plafond sans le franchir durablement) et
  un régime à croissance soutenue (qui franchit successivement les
  plafonds) ;
\item si la moyenne stationnaire augmente continûment jusqu'à croiser un
  plafond arbitraire, la cellule est classée « censurée par le garde-fou »
  et non « explosive ».
\end{enumerate}

Le rapport final montrera le diagnostic par des trajectoires et une carte
locale, et distinguera formellement changement de régime, forte population
stationnaire et censure numérique.

\section{Règles de passage à la confirmation et critères d'arrêt}
\label{sec:confirmation}

\subsection{Sélection des cellules confirmatoires}

Sont retenues pour la phase confirmatoire : (i) le centre et la référence ;
(ii) l'axe de taille $\lambda\in\{10,30,100\}$ ; (iii) les cellules aux
plus grands effets standardisés du screening ; (iv) des paires encadrant
chaque frontière de régime identifiée ; (v) les interactions candidates des
coupes 2D ; (vi) le contrôle $\sigma=0$. La liste précise, arrêtée après le
screening, sera consignée dans le rapport préliminaire avec sa
justification.

\subsection{Critères de confirmation}

Un effet exploratoire est dit \emph{confirmé} si, sur les graines
$\{11,\dots,15\}$ à $T=4000$ : son signe et son ordre de grandeur
persistent ; le contraste apparié moyen dépasse deux fois l'écart-type
inter-graines du contraste (ou l'intervalle à 95\,\% exclut zéro) ; il est
insensible au choix du burn-in ($T/8$, $T/4$, $T/2$) et cohérent entre
demi-fenêtres. Sinon, l'effet est rapporté comme non confirmé --- résultat
négatif à part entière.

\subsection{Critères d'arrêt de la campagne}

La campagne s'arrête lorsque : les volets 1--7 sont exécutés ou
explicitement comptabilisés comme échecs ou censures ; chaque censure
\code{pop\_max} a reçu son diagnostic ; les cellules confirmatoires ont
leurs cinq graines (ou plus près des frontières) ; les métriques primaires
du centre ont un écart-type inter-graines documenté ; et les rapports sont
compilés. Aucune cellule n'est exclue silencieusement : toute cellule
plannifiée apparaît dans les manifestes avec son statut final.

\section{Audit du moteur, benchmark et budget}
\label{sec:audit}

\subsection{Audit (exécuté le 17 juillet 2026)}

\begin{itemize}
\item \textbf{Tests moteur} : \code{tests/test\_mini.py}, 6 tests OK en
  10\,s, dont la \emph{parité exacte M4/M4B} sur deux régimes (graine 0,
  120 pas, $\lambda=10$, $\sigma=0{,}25$, $k=3$ ; graine 7, 90 pas,
  $\lambda=18$, $\sigma=0{,}31$, $k=2$) : séries, états individuels,
  contrats et avalanches identiques champ à champ. C'est l'unique recours
  exécutable à M4 de l'étude (contrôle de provenance).
\item \textbf{Tests statistiques} : \code{tests/test\_reporting.py},
  4 tests OK (manifeste de figures, histogramme adaptatif, Gini/Lorenz,
  MLE tronquée).
\item \textbf{Neutralité des mesures} : deux contrôles dédiés --- même
  configuration exécutée avec écriture minimale puis complète
  (\code{snapshot\_every} 0 vs 50, \code{individual\_every} 0 vs 1) :
  \code{series.csv} bit à bit identiques ($\lambda=10$ et $\lambda=30$,
  graine 0).
\item \textbf{Intégrité} : sur les six runs de benchmark, résidu comptable
  relatif maximal $\le 6\cdot10^{-16}$, zéro erreur de carnet.
\item \textbf{Provenance} : le moteur n'est pas suivi par git ; la
  provenance est fixée par \code{MODEL\_VERSION} (\code{m4b-mini-3}) et par
  le condensat SHA-256 des trois fichiers du moteur, tronqué à 16
  caractères : \code{f3a7ce02cd9185c7}. Ce condensat est inscrit dans
  chaque manifeste de run et vérifié avant toute réutilisation.
\item \textbf{Inventaire des runs existants} : Simulation Lab contient 10
  runs M4B, tous des validations d'intégration ($T\le260$, sauf un run de
  référence $T=2000$ à une seule graine) ; aucun ne satisfait les critères
  de réutilisation d'une cellule de campagne (panel de graines, phase,
  écriture contrôlée). La campagne ne réutilise donc aucun run existant ;
  tout est reproduit sous manifeste.
\end{itemize}

\subsection{Benchmark de coût (graine 0, $T=2000$, écriture minimale)}

\begin{center}
\begin{tabular}{@{}lrrr@{}}
\toprule
Cellule & Durée (s) & Population finale & Disque \\
\midrule
$\lambda=10$ centre micro. & 24 & 156 & 4{,}0 Mo \\
$\lambda=30$ & 80 & 441 & 11 Mo \\
$\lambda=100$ & 291 & 1\,569 & 36 Mo \\
$\lambda=30$, instantanés/100 pas & 83 & 441 & 13 Mo \\
$\lambda=10$, écriture complète & 39 & 156 & 30 Mo \\
$\lambda=10$, $\delta=0{,}02$, $\sigma=0{,}10$ & 42 & 230 & 4{,}1 Mo \\
\bottomrule
\end{tabular}
\end{center}

Le coût croît approximativement comme $\lambda\,T$ (population
proportionnelle à $\lambda$, marché en un round par tête). L'écriture
d'instantanés périodiques est quasi gratuite ; la table longitudinale
individuelle complète domine le disque et reste réservée aux besoins
explicites.

\subsection{Budget}

Machine : 8 cœurs, 31 Go de RAM, $\sim$106 Go libres. Au plus 6 processus
simultanés (2 cœurs laissés au système), un processus par run. Budget
prévisionnel de production : pilotes $\sim$30 runs ; OAT $\sim$70 runs ;
plan global 192 runs ; coupes 2D $\sim$200--300 runs ; confirmation
$\sim$60--80 runs ($T=4000$) ; diagnostics \code{pop\_max} selon besoin.
Soit de l'ordre de 600 runs, $\sim$15--20 h CPU, $\sim$2 à 3 h de mur en
parallèle, et $\sim$5--8 Go de disque en écriture minimale --- dans le
budget de la machine. Les runs de screening utilisent
\code{snapshot\_every=0} et \code{individual\_every=0} (neutralité
démontrée ci-dessus) ; les runs confirmatoires activent les instantanés
périodiques.

\section{Infrastructure et traçabilité}

\begin{itemize}
\item \textbf{Arborescence} : \path{scripts/} (exécution, métriques,
  analyses), \path{manifests/} (plans et échecs), \path{results/runs/}
  (données primaires par run), \path{results/metrics/} (JSON de
  métriques), \path{figures/} et \path{report/}.
\item \textbf{Identifiant de run} : \code{m4b\_} suivi des 12 premiers
  caractères hexadécimaux du SHA-256 de la configuration canonique
  complète (paramètres, graine, horizon, \code{pop\_max}, options
  d'écriture, condensat du moteur). Déterministe et non ambigu.
\item \textbf{Manifeste par run} : version et condensat du moteur,
  paramètres, graine, horizon, options d'écriture, statut du modèle,
  $t_{\mathrm{fin}}$, population finale, comptes globaux, erreurs de
  carnet, durée, taille disque, chemin des données, date et machine.
\item \textbf{Reprise} : les scripts d'exécution sautent toute cellule
  dont le manifeste est valide (même condensat moteur) et relancent le
  reste ; ils sont relançables après interruption.
\item \textbf{Journal technique} : \code{JOURNAL.md} dans le dossier de
  campagne, tenu au fil de l'eau (décisions, anomalies, résultats
  négatifs).
\item \textbf{Livrables} : les interprétations scientifiques sont dans
  les trois rapports LaTeX (protocole, rapport préliminaire, rapport
  final) ; les CSV/JSON sont des annexes techniques. Chaque
  figure et tableau des rapports doit être reproductible depuis les
  manifestes ; l'annexe de traçabilité du rapport final liste chaque run
  cité avec son identifiant, ses paramètres, sa graine et son statut.
\end{itemize}

\section*{Hors périmètre : horizon M5}
\addcontentsline{toc}{section}{Hors périmètre : horizon M5}

L'étude ne modifie pas le moteur, n'implémente aucune réduction de
paramètres, ne fait varier aucun paramètre en cours de run et n'annonce
aucun effet rebond. Le rapport final devra en revanche livrer : les
baselines et unités qui permettront d'isoler l'effet du futur exposant de
concavité $\gamma$ ($F_\gamma(K)=\alpha K^\gamma$, $\gamma=1/2$ étant le
cas M4B) ; les variables nécessaires à une future mesure d'effet rebond au
sens de Sorrell et Dimitropoulos (2008) ; et une feuille de route pour
l'étude analytique (points fixes, échelles, conditions de stationnarité,
prédictions falsifiables).

\end{document}
